\documentclass[11pt]{amsart}

\usepackage[T1]{fontenc}
\usepackage{lmodern}
\usepackage{microtype}
\usepackage{mathtools,amssymb,amsthm}
\usepackage[hidelinks]{hyperref}

\hypersetup{
  pdftitle={The Kiani-Tavakolipour Dichotomy at m = 5 for Symmetric Tropical Matrices, in Additive Form}
}

\newtheorem{theorem}{Theorem}
\newtheorem{conjecture}[theorem]{Conjecture}
\newtheorem{lemma}[theorem]{Lemma}

\newcommand{\Rmax}{\mathbb R_{\max}}
\newcommand{\dd}{\delta}
\newcommand{\Cov}{\mathcal C}
\newcommand{\Ee}{E}
\newcommand{\Allowed}{\mathcal A}
\newcommand{\per}{\operatorname{per}}
\newcommand{\supp}{\operatorname{supp}}
\newcommand{\pos}{\operatorname{pos}}

\title[The Kiani--Tavakolipour dichotomy at $m=5$, in additive form]
{The Kiani--Tavakolipour Dichotomy at $m=5$ for Symmetric Tropical Matrices,
in Additive Form}

\date{}

\subjclass[2020]{15A80, 14T90, 90C05}
\keywords{tropical characteristic polynomial, max-plus algebra, symmetric matrix,
principal submatrix, tropical permanent, concavity, Gordan's theorem}

\begin{document}

\begin{abstract}
For an $n\times n$ matrix $A$ over $\Rmax=\mathbb R\cup\{-\infty\}$ let
$\dd_k(A)$ be the largest tropical permanent of a principal $k\times k$
submatrix of $A$, with $\dd_0=0$; these are the coefficients of the tropical
characteristic polynomial. Kiani and Tavakolipour conjecture that for symmetric
$A$ and every $m$ with $3\le m\le n$ at least one of
$\dd_m-\dd_{m-1}\le\dd_{m-1}-\dd_{m-2}$ and
$\dd_m-\dd_{m-2}\le2(\dd_{m-2}-\dd_{m-3})$ holds; they prove the cases $m=3$ and
$m=4$ and leave $m\ge5$ open. We prove that for every $n\ge5$ and every
symmetric $A$ over $\Rmax$, at least one of the additive inequalities
$\dd_5+\dd_3\le2\dd_4$ and $\dd_5+2\dd_2\le3\dd_3$ holds. Whenever the $\dd_k$
involved are finite these are equivalent to the two conjectured inequalities at
$m=5$, so the conjecture is settled at $m=5$ in every case in which its
differences are
defined; when some $\dd_k=-\infty$ the additive inequalities are our chosen
subtraction-free extension of the source statement rather than that statement
itself. The failure of both additive inequalities is first rewritten as a system
of strict homogeneous linear inequalities in the entries of $A$; a nonnegative
rational combination of those forms that vanishes identically then refutes the
failure. A failure is witnessed by three optimal selections of sizes $5$, $3$
and $2$, whose union meets at most $10$ indices, and entrywise nonnegativity
shows that a certificate may be sought inside that index set, which handles all
$n$ at once. What remains is a finite check:
$7\times600\times90=378{,}000$ triples, falling into $53{,}900$ classes with
identical constraint systems, each carrying such a combination; $38{,}524$ of
the classes are killed by a two-term identity between multisets of matrix
positions, the rest by longer identities or a general rational combination. The
cases $m\ge6$ remain open.
\end{abstract}

\maketitle

\section{The conjecture}

Let $\Rmax=\mathbb R\cup\{-\infty\}$ with $a\oplus b=\max(a,b)$ and
$a\odot b=a+b$. For a $k\times k$ matrix $B$ over $\Rmax$ the tropical permanent
is $\per(B)=\max_{\pi\in S_k}\sum_{i=1}^k b_{i\pi(i)}$. For an $n\times n$
matrix $A$ over $\Rmax$ and $0\le k\le n$ set
\begin{equation}\label{eq:delta}
 \dd_k(A)=\max_{S\subseteq[n],\,|S|=k}\per\bigl(A[S]\bigr),
 \qquad \dd_0(A)=0,
\end{equation}
where $A[S]$ is the principal submatrix on the rows and columns indexed by $S$.
Up to sign these are the coefficients of the tropical characteristic polynomial
of $A$; \eqref{eq:delta} is the description used throughout
\cite[Thm.~2.6 and eq.~(2.3)]{KT}. We write $\dd_k$ for $\dd_k(A)$.

The conjecture below is Conjecture~6.1 of \cite{KT}, stated there for ``the
sequence $(\dd_m)$ associated with \emph{symmetric} tropical matrices''. It
reads, verbatim:

\begin{conjecture}[Kiani--Tavakolipour {\cite[Conj.~6.1]{KT}}]\label{conj:KT}
For $3\le m\le n$, at least one of the following inequalities holds:
\begin{align}
 \dd_m-\dd_{m-1}&\le\dd_{m-1}-\dd_{m-2},\quad\text{or}\label{eq:C1}\\
 \dd_m-\dd_{m-2}&\le2(\dd_{m-2}-\dd_{m-3}).\label{eq:C2}
\end{align}
\end{conjecture}

\medskip
\noindent\textbf{The additive reading.}
The semiring is $\Rmax$, so a matrix may have $-\infty$ entries and a $\dd_k$
may equal $-\infty$; in that case the differences in \eqref{eq:C1} and
\eqref{eq:C2} have no literal meaning, and we make no claim about what
Conjecture~\ref{conj:KT} asserts there. We work instead with the following
additive forms, obtained by moving the negative terms to the other side. These
are the statements we prove:
\begin{align}
 \dd_m+\dd_{m-2}&\le2\dd_{m-1},\label{eq:I}\tag{I}\\
 \dd_m+2\dd_{m-3}&\le3\dd_{m-2}.\label{eq:II}\tag{II}
\end{align}
Whenever all the $\dd_k$ involved are finite, \eqref{eq:I} is equivalent to
\eqref{eq:C1} and \eqref{eq:II} to \eqref{eq:C2}; outside that case
\eqref{eq:I}--\eqref{eq:II} are our chosen extension of the source's
subtraction-based statement, not a reading we attribute to \cite{KT}.

Conjecture~\ref{conj:KT} is proved in \cite{KT} for $m=3$ (their Thm.~4.1) and
$m=4$ (their Thm.~4.5); $m=5$ is the first case the source leaves open.

\begin{theorem}\label{thm:main}
Let $n\ge5$ and let $A$ be a symmetric $n\times n$ matrix over $\Rmax$. Then
\begin{equation*}
 \dd_5+\dd_3\le2\dd_4
 \qquad\text{or}\qquad
 \dd_5+2\dd_2\le3\dd_3 .
\end{equation*}
\end{theorem}

\noindent That is, at least one of \eqref{eq:I}, \eqref{eq:II} holds at $m=5$,
for every $n$ and with no finiteness assumption on $A$. Whenever
$\dd_5,\dd_4,\dd_3,\dd_2$ are finite this is exactly the assertion of
Conjecture~\ref{conj:KT} at $m=5$, so the conjecture is settled at $m=5$ in
every case in which its differences are defined; in the remaining cases the
theorem asserts the additive reading \eqref{eq:I}--\eqref{eq:II} and not the
source statement.

\section{Covers}

A \emph{$k$-cover} of $[n]$ is a pair $P=(S,\pi)$ with $S\subseteq[n]$,
$|S|=k$, and $\pi$ a permutation of $S$; write $\supp P=S$ and
$w(P)=\sum_{i\in S}a_{i\pi(i)}\in\Rmax$. Let $\Cov_k$ be the set of $k$-covers.
Reading \eqref{eq:delta} through the permanent gives
\begin{equation}\label{eq:cover}
 \dd_k=\max_{P\in\Cov_k}w(P),
\end{equation}
the point being that a principal submatrix forces the row set and the column set
to coincide. The unique $0$-cover has weight $0$, consistent with $\dd_0=0$.

Since $A$ is symmetric, $w(P)$ depends on $P$ only through the multiset
\[
 \Ee(P)=\bigl\{\!\{\,\{i,\pi(i)\}\;:\;i\in\supp P\,\}\!\bigr\}
\]
of \emph{positions} $\{i,j\}$ (with $i=j$ allowed), which we regard as a vector
of nonnegative integers indexed by positions; then
$w(P)=\sum_p\Ee(P)_p\,a_p$. A loop $i\mapsto i$ contributes $\{i,i\}$ once; a
transposition $(i\,j)$ contributes $\{i,j\}$ twice, and indeed weight
$2a_{ij}$. Write $\pos P$ for the \emph{set} of positions occurring in
$\Ee(P)$. A $k$-cycle and its inverse have the same $\Ee$, so for symmetric $A$
the relevant object is $\Ee(P)$, not $P$.

\section{The failure of both inequalities, as a linear system}

Throughout this section $m=5$; nothing in the argument is special to $5$.

\begin{lemma}\label{lem:reform}
Both \eqref{eq:I} and \eqref{eq:II} fail at a symmetric $A$ if and only if there
are $P\in\Cov_5$, $Q\in\Cov_3$, $S\in\Cov_2$ with
\begin{align}
 w(P)+w(Q)&>2w(R) &&\text{for every }R\in\Cov_4,\label{eq:sysI}\\
 w(P)+2w(S)&>3w(Q') &&\text{for every }Q'\in\Cov_3.\label{eq:sysII}
\end{align}
Moreover $w(P),w(Q),w(S)$ are then finite.
\end{lemma}

\begin{proof}
Suppose both fail, i.e.\ $\dd_5+\dd_3>2\dd_4$ and $\dd_5+2\dd_2>3\dd_3$. In
$\Rmax$ a left-hand side equal to $-\infty$ exceeds nothing, so
$\dd_5,\dd_3,\dd_2$ are finite. Take $P,Q,S$ attaining the maxima in
\eqref{eq:cover} for $k=5,3,2$; their weights are $\dd_5,\dd_3,\dd_2$, finite.
For $R\in\Cov_4$, $2w(R)\le2\dd_4<\dd_5+\dd_3=w(P)+w(Q)$, which is
\eqref{eq:sysI}; \eqref{eq:sysII} is identical with $\dd_3,\dd_2$ in place of
$\dd_4,\dd_3$.

Conversely, given such $P,Q,S$, the right-hand sides of \eqref{eq:sysI} and
\eqref{eq:sysII} already range over everything the maxima $\dd_4,\dd_3$ range
over, so $2\dd_4<w(P)+w(Q)\le\dd_5+\dd_3$ and
$3\dd_3<w(P)+2w(S)\le\dd_5+2\dd_2$, using $\dd_5\ge w(P)$, $\dd_3\ge w(Q)$,
$\dd_2\ge w(S)$ and monotonicity of $\oplus$-addition. Finiteness of $w(P)$,
$w(Q)$, $w(S)$ is forced by \eqref{eq:sysI} and \eqref{eq:sysII} themselves.
\end{proof}

Fix such a triple and set
\[
 \Allowed=\pos P\cup\pos Q\cup\pos S,
 \qquad D=|\Allowed|\le5+3+2=10 .
\]
Because $w(P),w(Q),w(S)$ are finite, \emph{every entry $a_p$ with
$p\in\Allowed$ is finite}. Put
\[
 \mathrm{Pool}_4=\{R\in\Cov_4:\pos R\subseteq\Allowed\},
 \qquad
 \mathrm{Pool}_3=\{Q'\in\Cov_3:\pos Q'\subseteq\Allowed\},
\]
where multiplicities are unconstrained: a transposition on $\{i,j\}$ uses the
single position $\{i,j\}$ twice and belongs to the pool as soon as
$\{i,j\}\in\Allowed$. For $R\in\mathrm{Pool}_4$ and $Q'\in\mathrm{Pool}_3$
define integer vectors in $\mathbb Z^{\Allowed}$,
\begin{equation}\label{eq:forms}
 f_R=\Ee(P)+\Ee(Q)-2\Ee(R),
 \qquad
 g_{Q'}=\Ee(P)+2\Ee(S)-3\Ee(Q').
\end{equation}
Evaluated at the vector $x=(a_p)_{p\in\Allowed}\in\mathbb R^D$ of \emph{finite}
entries, $f_R(x)=w(P)+w(Q)-2w(R)$ and $g_{Q'}(x)=w(P)+2w(S)-3w(Q')$; so
\eqref{eq:sysI} and \eqref{eq:sysII}, restricted to the pool, say exactly that
$x$ makes every form in \eqref{eq:forms} strictly positive.

\begin{lemma}[certificates]\label{lem:cert}
Let $\Ee(P),\Ee(Q),\Ee(S)$ be given, with $\Allowed$ and the pools as above.
Suppose there are nonnegative rationals $\lambda_R$ $(R\in\mathrm{Pool}_4)$ and
$\mu_{Q'}$ $(Q'\in\mathrm{Pool}_3)$, not all zero, with
\begin{equation}\label{eq:certificate}
 \sum_R\lambda_Rf_R+\sum_{Q'}\mu_{Q'}g_{Q'}=0
 \qquad\text{in }\mathbb Q^{\Allowed}.
\end{equation}
Then no symmetric $A$, over any index set containing $\Allowed$, satisfies
\eqref{eq:sysI} and \eqref{eq:sysII} for this triple.
\end{lemma}

\begin{proof}
If it did, all entries indexed by $\Allowed$ would be finite by
Lemma~\ref{lem:reform} and each form would take a strictly positive real value
at $x$; evaluating \eqref{eq:certificate} at $x$ gives
$0=\sum\lambda_Rf_R(x)+\sum\mu_{Q'}g_{Q'}(x)>0$, since the coefficients are
nonnegative and not all zero. Contradiction.
\end{proof}

Normalising $\sum\lambda_R+\sum\mu_{Q'}=1$ and writing
$L=\sum\lambda_R$, $M=\sum\mu_{Q'}$, \eqref{eq:certificate} is the identity of
position multisets
\begin{equation}\label{eq:identity}
 \Ee(P)+L\,\Ee(Q)+2M\,\Ee(S)
 = 2\sum_R\lambda_R\Ee(R)+3\sum_{Q'}\mu_{Q'}\Ee(Q') ,
\end{equation}
with $L+M=1$, $L,M\ge0$. Suppose $A$ is symmetric and $P,Q,S$ \emph{attain} the
maxima in \eqref{eq:cover} for $k=5,3,2$, so that $w(P)=\dd_5$, $w(Q)=\dd_3$ and
$w(S)=\dd_2$, all finite, as for the triple supplied by
Lemma~\ref{lem:reform}. Evaluating \eqref{eq:identity} at $A$ and using
$\dd_4\ge w(R)$ for $R\in\mathrm{Pool}_4$, $\dd_3\ge w(Q')$ for
$Q'\in\mathrm{Pool}_3$, together with $\sum_R\lambda_R=L$,
$\sum_{Q'}\mu_{Q'}=M$ and $L+M=1$, then gives
$L(\dd_5+\dd_3-2\dd_4)+M(\dd_5+2\dd_2-3\dd_3)\le0$, so one of \eqref{eq:I},
\eqref{eq:II} holds. Checking one certificate requires only that the
multipliers be nonnegative and that a rational linear identity hold; but the
multipliers are printed below only for the sample certificates, so for the rest
of the census that check is available only by running the program. No
linear-programming
duality is needed: only the trivial direction of Gordan's theorem
\cite{Gordan} is used, namely that \eqref{eq:certificate} obstructs a strictly
positive solution; the converse is not used.

\section{All $n$ at once: the window $[3m-5]$}

\begin{lemma}\label{lem:window}
Suppose both \eqref{eq:I} and \eqref{eq:II} fail at some symmetric $A$ over
$\Rmax$ of some size $n\ge5$. Then there are $P_0\in\Cov_5([5])$,
$Q_0\in\Cov_3([10])$, $S_0\in\Cov_2([10])$, with $P_0$ a representative of its
cycle type, such that the triple $(\Ee(P_0),\Ee(Q_0),\Ee(S_0))$ admits no
certificate \eqref{eq:certificate}.
\end{lemma}

\begin{proof}
Take $P,Q,S$ from Lemma~\ref{lem:reform}. Their supports have union of size at
most $5+3+2=10$, and $|\supp P|=5$. Choose a bijection of that union with a
subset of $[10]$ carrying $\supp P$ onto $[5]$, then compose with a permutation
of $[5]$, extended by the identity on $\{6,\dots,10\}$, carrying $P$ to the
chosen representative of its cycle type. Relabelling the index set carries
$\Allowed$, $\mathrm{Pool}_4$ and $\mathrm{Pool}_3$ bijectively onto the
relabelled ones and permutes the coordinates of every form in
\eqref{eq:forms}, so a certificate for the image triple would be a certificate
for the original, contradicting Lemma~\ref{lem:cert} together with
\eqref{eq:sysI}--\eqref{eq:sysII}. The images $Q_0,S_0$ are covers of $[10]$ of
sizes $3$ and $2$.
\end{proof}

The window comes from the union of the supports of $P$, $Q$ and $S$:
$|\Allowed|\le m+(m-2)+(m-3)=3m-5$, which is $10$ at $m=5$. Entrywise
nonnegativity is what makes it enough to search there. In the identity
\eqref{eq:identity} every vector has nonnegative entries, so at a position
outside $\Allowed$ the left side is $0$ and therefore every $R$ with
$\lambda_R>0$ and every $Q'$ with $\mu_{Q'}>0$ vanishes there; a certificate
never leaves the window, so restricting the pools to it loses nothing.

\section{The finite check}

By Lemma~\ref{lem:window}, Theorem~\ref{thm:main} follows once a certificate is
exhibited for every triple $(\Ee(P_0),\Ee(Q_0),\Ee(S_0))$ in the following
finite list. The list is completely specified by the three counts below, and a
reader can regenerate it without any code.

\medskip
\noindent\emph{$P_0$: seven triples' worth of cycle type.}
Two $5$-covers of a $5$-set have relabelling-equivalent position multisets
precisely when they have the same cycle type, since a cycle and its inverse have
the same $\Ee$ and relabelling preserves cycle type. So $P_0$ runs over the
\textbf{seven} partitions of $5$,
\[
 1^5,\quad 2\cdot1^3,\quad 2^2\cdot1,\quad 3\cdot1^2,\quad 3\cdot2,
 \quad 4\cdot1,\quad 5 .
\]

\medskip
\noindent\emph{$Q_0$: the $\mathbf{600}$ distinct $\Ee$ of a $3$-cover of
$[10]$.}
Three loops, $\binom{10}{3}=120$; a transposition and a loop,
$\binom{10}{2}\cdot8=360$; a $3$-cycle, one position multiset per $3$-subset
because a $3$-cycle and its inverse agree, $\binom{10}{3}=120$. Total
$120+360+120=600$.

\medskip
\noindent\emph{$S_0$: the $\mathbf{90}$ distinct $\Ee$ of a $2$-cover of
$[10]$.}
Two loops, $\binom{10}{2}=45$; a transposition, $45$. Total $90$.

\medskip
\noindent The list therefore has
\[
 7\times600\times90=378{,}000
\]
entries. Every one of them carries a certificate \eqref{eq:certificate}; none is
live. This is the computational content of the paper and it is what
\texttt{verify.py} in the accompanying folder re-derives from scratch. Under the
order-preserving relabelling of the index set the $378{,}000$ triples fall into
$53{,}900$ classes with identical constraint systems, and the program solves one
system per class.

The certificates are small and mostly of two shapes. Writing $ij$ for the
position $\{i,j\}$:

\begin{itemize}
\item $38{,}524$ of the $53{,}900$ classes are killed by a
\emph{two-term identity in branch \textup{(I)}}: two pool $4$-covers $R_1,R_2$
with
\[
 \Ee(P)+\Ee(Q)=\Ee(R_1)+\Ee(R_2),
\]
i.e.\ \eqref{eq:certificate} with $\lambda_{R_1}=\lambda_{R_2}=\tfrac12$,
$L=1$, $M=0$.
\item $1{,}016$ classes are killed by a \emph{three-term identity in branch
\textup{(II)}}: pool $3$-covers with
$\Ee(P)+2\Ee(S)=\Ee(Q'_1)+\Ee(Q'_2)+\Ee(Q'_3)$, i.e.\
$\mu_{Q'_i}=\tfrac13$, $L=0$, $M=1$.
\item the remaining $14{,}360$ classes need a general nonnegative rational
combination. Over all $53{,}900$ classes the largest denominator occurring in
any multiplier is $48$, so every certificate is a small rational vector.
\end{itemize}

Two certificates in full, hand-checkable by counting symbols. In each, $P$ is
the all-loops $5$-cover of $\{0,1,2,3,4\}$.
\begin{align*}
 \Ee(P)+\Ee(Q)
 &=\{00,11,22,33,44\}+\{00,11,22\}\\
 &=\{00,00,11,11,22,22,33,44\}\\
 &=\{00,11,22,33\}+\{00,11,22,44\}
 =\Ee(R_1)+\Ee(R_2),
\end{align*}
with $Q$ the three loops on $\{0,1,2\}$ and $R_1,R_2$ the all-loops $4$-covers
of $\{0,1,2,3\}$ and $\{0,1,2,4\}$; and, with $Q$ the transposition $(0\,1)$
together with the loop at $2$,
\begin{align*}
 \Ee(P)+\Ee(Q)
 &=\{00,11,22,33,44\}+\{01,01,22\}\\
 &=\{01,01,22,33\}+\{00,11,22,44\}
 =\Ee(R_1)+\Ee(R_2),
\end{align*}
where now $R_1$ is the transposition $(0\,1)$ with loops at $2,3$. In both
cases $L=1$, $M=0$ and the two multipliers are $\tfrac12$, so
\eqref{eq:identity} gives $\dd_5+\dd_3\le2\dd_4$ at $A$; that is, branch~(I)
holds.

\begin{proof}[Proof of Theorem~\ref{thm:main}]
If the conclusion failed for some symmetric $A$ and some $n\ge5$, then by
Lemma~\ref{lem:window} some triple in the list above would admit no certificate
\eqref{eq:certificate}. Every triple in the list does admit one. Hence no such
$A$ exists.
\end{proof}

\section{What is not settled}

\noindent\textbf{$m\ge6$.} Theorem~\ref{thm:main} settles the single value
$m=5$. Conjecture~\ref{conj:KT} ranges over $3\le m\le n$ and remains open for
every $m\ge6$. The reduction above is generic in $m$ and gives the window
$[3m-5]=[13]$ at $m=6$, with a raw count of
$11\times|\Cov_4([13])|\times|\Cov_3([13])|=11\times12{,}155\times1{,}430
\approx1.9\times10^8$ triples, which we have not attempted. Nothing here
suggests the answer at $m\ge6$ either way.

\medskip
\noindent\textbf{Nonsymmetric matrices.} Symmetry is essential and \cite{KT}
says so; Theorem~\ref{thm:main} makes no claim about nonsymmetric $A$.

\medskip
\noindent\textbf{The multiplier vectors are not shipped.}
Only the two certificates above appear in print. The multipliers for the rest of
the $53{,}900$ classes exist only inside a run of \texttt{verify.py}, which
regenerates and re-verifies them but writes no file listing them, so the census
is reproducible without being auditable from data in this folder.

\medskip
\noindent\textbf{Only one direction of Gordan's theorem is used.} We never claim
that a triple \emph{without} a certificate yields a genuine counterexample.

\section{Relation to earlier work}

Lewis \cite{Lewis} studies the same object under the name \emph{best principal
submatrix problem}, settling the symmetric case with entries in
$\{0,-\infty\}$ and the symmetric case whose digraph has no odd cycle; neither
states an inequality of the shape \eqref{eq:C1}--\eqref{eq:C2}. Rosenmann
\cite{Rosenmann} proves concavity in $k$ for the bipartite $k$-cardinality
assignment problem, where the $k$ rows and the $k$ columns are chosen
independently, whereas \eqref{eq:delta} forces the two sets to coincide.

\section{Reproduction and verification}

Sections~2--5 are self-contained: the census is specified by the seven cycle
types, $\Cov_3([10])$ and $\Cov_2([10])$, and a reader can regenerate it in a
few lines and search for the certificates directly. The accompanying
\texttt{verify.py} does exactly that, in Python~3 with the standard library only
and in exact integer and \texttt{Fraction} arithmetic. It re-derives the counts
$7$, $600$, $90$ and $378{,}000$; checks the printed certificates as exact
multiset identities; and re-runs the $m=5$ census, re-substituting every one of
the $53{,}900$ multiplier vectors it finds into its own forms. It closes with
\texttt{VERDICT: ALL 47 CHECKS PASS}; \texttt{verify.output.txt} in the same
folder records one such run.

The program prints its own limits: it does not touch $m\ge6$, it does not
persist the multiplier vectors it verifies, and it makes no literature claim.

\begin{thebibliography}{9}

\bibitem{Gordan}
P.~Gordan,
\emph{\"Uber die Aufl\"osung linearer Gleichungen mit reellen Coefficienten},
Math. Ann. \textbf{6} (1873), 23--28.
\href{https://doi.org/10.1007/BF01442864}{doi:10.1007/BF01442864}.

\bibitem{KT}
D.~Kiani and H.~Tavakolipour,
\emph{Properties of the tropical characteristic polynomial of symmetric
matrices},
arXiv:2607.10922v1, 12 July 2026.
\href{https://arxiv.org/abs/2607.10922}{arXiv:2607.10922}.
Statement numbers are those of that e-print, the only text we could access.

\bibitem{Lewis}
S.~C. Lewis,
\emph{On the best principal submatrix problem},
Ph.D. thesis, University of Birmingham, 2007.
\href{https://etheses.bham.ac.uk/id/eprint/26}{etheses.bham.ac.uk/id/eprint/26}.

\bibitem{Rosenmann}
A.~Rosenmann,
\emph{Computing the sequence of $k$-cardinality assignments},
J. Comb. Optim. \textbf{44} (2022), 1265--1283.
\href{https://doi.org/10.1007/s10878-022-00840-7}{doi:10.1007/s10878-022-00840-7};
\href{https://arxiv.org/abs/2104.04037}{arXiv:2104.04037}.

\end{thebibliography}

\end{document}
